% Generated by scripts/assemble.py; corrections belong in source/corrections.json.
\examhead{146}{中南大学 2026 年数学分析试卷(回忆版)}{139}
% p140-b01; source PDF page 140
\blockanchor{p140-b01}
(注. 题目来自考生回忆, 部分表述或数据可能有误)\par
% p140-b02; source PDF page 140
\par\Needspace{5\baselineskip}
\blockanchor{p140-b02}
\majorquestion{一}计算题. 每题 10 分, 共 40 分.\par
% p140-b03; source PDF page 140
\blockanchor{p140-b03}
\question{1}设 $x = 1 + t^3$, $y = \ln (1 + t)$, 求 $\left.\frac{\mathrm{d}y}{\mathrm{d}x}\right|_{t=1}$.\par
% p140-b04; source PDF page 140
\blockanchor{p140-b04}
\question{2}求极限 $\lim_{x\to 0}\frac{\int_x^{x^3}\ln(1 + t^2)\mathrm{d}t}{x^3}.$\par
% p140-b05; source PDF page 140
\blockanchor{p140-b05}
\question{3}求积分 $\int_0^{+\infty} \frac{2x \ln(1 + x^2)}{(1 + x^2)^2} \mathrm{d}x$.\par
% p140-b06; source PDF page 140
\blockanchor{p140-b06}
\question{4}求级数 $\sum_{n=1}^{\infty} \frac{1}{n(n+2)}$ 的和.\par
% p140-b07; source PDF page 140
\par\Needspace{5\baselineskip}
\blockanchor{p140-b07}
\majorquestion{二}(20 分) 设 $f: [0,1]^3 \to \mathbb{R}$ 为连续函数, 证明: $f$ 的像为有界闭区间.\par
% p140-b08; source PDF page 140
\par\Needspace{5\baselineskip}
\blockanchor{p140-b08}
\majorquestion{三}(20 分) 设 $f(x)$ 在 $[0,1]$ 上连续, 在 $(0,1)$ 上二阶可导, $f(0) = 0$, $f\left(\frac{1}{2}\right) = \frac{1}{2}$, $f(1) = 0$, 证明: 存在 $\xi \in (0,1)$, 满足 $f''(\xi) = -4$.\par
% p140-b09; source PDF page 140
\par\Needspace{5\baselineskip}
\par\Needspace{6\baselineskip}
\blockanchor{p140-b09}
\majorquestion{四}(20 分) 求积分\par
\nopagebreak[4]
% p140-b10; source PDF page 140
\blockanchor{p140-b10}
\displayformula{\iint_ {\Sigma} x \mathrm{d} y \mathrm{d} z + y \mathrm{d} z \mathrm{d} x + z \mathrm{d} x \mathrm{d} y.}
% p140-b11; source PDF page 140
\blockanchor{p140-b11}
其中  $\Sigma$  为球面  $x^{2} + y^{2} + z^{2} = 25$ ，取曲面的内侧.\par
% p140-b12; source PDF page 140
\par\Needspace{5\baselineskip}
\blockanchor{p140-b12}
\majorquestion{五}(20 分) 设 $f(x)$ 在 $[1, +\infty)$ 上一致连续, $\varphi(x)$ 在 $[1, +\infty)$ 上连续, 且 $\lim_{x \to +\infty} [f(x) - \varphi(x)] = 0$, 证明: $\varphi(x)$ 在 $[1, +\infty)$ 上一致连续.\par
% p140-b13; source PDF page 140
\par\Needspace{5\baselineskip}
\blockanchor{p140-b13}
\majorquestion{六}(10 分) 求反常积分 $\int_0^{+\infty} \frac{e^{-ax} \sin bx}{x} \, \mathrm{d}x, a > 0, b \in \mathbb{R}$.\par
% p140-b14; source PDF page 140
\par\Needspace{5\baselineskip}
\blockanchor{p140-b14}
\majorquestion{七}(10 分) 设 $S(a) = \sum_{n=1}^{\infty} \frac{\sin^2(n\pi a)}{n^3}, a \in [0,1]$.\par
% p140-b15; source PDF page 140
\blockanchor{p140-b15}
(1) 求 $S(a)$ 的收敛域, 并证明在收敛域内 $S(a)$ 一致收敛.\par
% p140-b16; source PDF page 140
\blockanchor{p140-b16}
(2) 证明: $S(a)$ 在 $[0,1]$ 上连续可微.\par
% p140-b17; source PDF page 140
\blockanchor{p140-b17}
(3) 证明: $f(a) = \frac{S(a)}{a^2}$ 单调递减. 提示: $S''(a) = \sum_{n=1}^{\infty} \frac{2\pi^2 \cos(2n\pi a)}{n} = -\pi^2 \ln(4\sin^2(\pi a)).$\par
% p140-b18; source PDF page 140
\par\Needspace{5\baselineskip}
\blockanchor{p140-b18}
\majorquestion{八}(10 分) 已知在 $\mathbb{R}^3$ 中, 在点 $x_0$ 处存在某个电子, 由该电子产生的电场为 $E = \frac{q}{4\pi\varepsilon_0}\frac{x - x_0}{\|x - x_0\|^3}$, 其中 $x \in \mathbb{R}^3$, 且 $x \neq x_0, \varepsilon_0$ 是正常数. 设 $\Omega$ 为 $\mathbb{R}^3$ 中的单连通区域, 其边界 $\partial \Omega$ 为光滑曲面.\par
% p140-b19; source PDF page 140
\blockanchor{p140-b19}
(1) 求证: $x_0 \notin \Omega$ 时, 电通量 $\Phi = \iint_{\partial \Omega} E \, \mathrm{d}S = 0$.\par
% p140-b20; source PDF page 140
\blockanchor{p140-b20}
(2) $x_0 \in \Omega$ 时, 电通量 $\Phi = \iint_{\partial \Omega} E \, \mathrm{d}S$ 为常数, 求此常数.\par
